\section{Clasificación de los ciclos aritméticos irreducibles de la red radical nonádica}

En la realizaci\'on aritm\'etica, un primo corresponde a un ciclo
irreducible que no se factoriza en ciclos positivos menores. La periodicidad
non\'adica puede seleccionar candidatos y organizar clases residuales; la
primalidad permanece definida por la propiedad multiplicativa exacta

\[
p=ab\Longrightarrow a=1\ \text{o}\ b=1.
\]

Una vez seleccionado el coproducto de ciclos aritm\'eticos primos, se
obtienen \cref{eq:meissel-mertens,eq:twin-prime-constant}. La estructura
helicoidal y la ra\'iz digital parametrizan fase y acarreo, pero no sustituyen
el criterio de factorizaci\'on.

\section{Densidad, t\'ermino regularizado y holonom\'ia de vacancias}

La densidad positiva b\'asica es

\begin{equation}
\varepsilon_{\rm vac}=\log_{10}\frac{10}{9}.
\label{eq:vacancy-density}
\end{equation}

La secuencia se deduce de la completaci\'on
\(1000=729+271\) determina
\begin{equation}
\varepsilon_{\rm vac}
=1-\log_{1000}(729)
=\log_{10}\frac{10}{9},
\label{eq:vacancy-density-two-forms}
\end{equation}
y la palabra mec\'anica exacta
\begin{equation}
z_n=\lfloor(n+1)\varepsilon_{\rm vac}\rfloor
-\lfloor n\varepsilon_{\rm vac}\rfloor,
\qquad n\ge1.
\label{eq:vacancy-mechanical-word}
\end{equation}
Como \(\varepsilon_{\rm vac}\) es irracional, la palabra es Sturmiana:
tiene complejidad \(p(N)=N+1\), balance uno y entrop\'ia topol\'ogica cero.
Esta din\'amica difiere del desplazamiento completo sobre el alfabeto
tri\'adico: la vacancia determina una frontera ordenada de entrop\'ia nula.

La suma telesc\'opica produce
\begin{equation}
\sum_{n=1}^{N}z_n
=\lfloor(N+1)\varepsilon_{\rm vac}\rfloor,
\qquad
\sum_{n=1}^{N}(z_n-\varepsilon_{\rm vac})
=\varepsilon_{\rm vac}
-\{(N+1)\varepsilon_{\rm vac}\}.
\label{eq:vacancy-discrepancy}
\end{equation}
Las sumas parciales de la discrepancia tienen m\'odulo menor que uno. Esta
es la estimaci\'on que sostiene las constantes arm\'onicas posteriores.

La fracci\'on continua comienza
\[
\varepsilon_{\rm vac}
=[0;21,1,5,1,6,2,3,\ldots].
\]
Por ello los unos de \(z_n\) est\'an separados por \(21\) o \(22\) pasos.
Si
\[
\sigma=22-\varepsilon_{\rm vac}^{-1}
=1-T_G(\varepsilon_{\rm vac}),
\]
entonces \(a_1(\sigma)=6\) y los retornos de la separaci\'on corta son de
longitud \(6\) o \(7\). Las parejas \(21/22\) y \(6/7\) son dos escalas de
una misma rotaci\'on, no coincidencias de enteros.

Para la palabra exacta \cref{eq:vacancy-mechanical-word}, la constante regularizada es

\begin{equation}
\gamma_{\rm vac}^{+}
=\lim_{N\to\infty}\left(
\sum_{n\le N}\frac{z_n}{n}
-\varepsilon_{\rm vac}\log N
\right).
\label{eq:vacancy-gamma}
\end{equation}

La holonom\'ia asociada es

\begin{equation}
\rho_\varepsilon(k)=\exp(2\pi\ii k\varepsilon_{\rm vac}).
\label{eq:vacancy-holonomy}
\end{equation}

Las tres cantidades tienen tipos distintos: densidad, t\'ermino finito y fase. La torre de Stieltjes de vacancias prolonga \cref{eq:vacancy-gamma}:

\begin{equation}
\gamma_{\rm vac,j}^{+}
=\varepsilon_{\rm vac}\gamma_j
+\sum_{n\ge1}\frac{(z_n-\varepsilon_{\rm vac})(\log n)^j}{n}.
\label{eq:vacancy-stieltjes}
\end{equation}

Para \(j=0\) se recupera \(\gamma_{\rm vac}^{+}\). La convergencia requiere la cota de discrepancia del desarrollo integral de vacancias.

La convergencia se prueba aqu\'i. Para \(j\ge0\), la sucesi\'on
\(b_n=(\log n)^j/n\) tiende a cero y es decreciente a partir de
\(n>e^j\). La cota uniforme \cref{eq:vacancy-discrepancy} y el criterio de
Dirichlet prueban la convergencia de
\[
\sum_{n\ge1}(z_n-\varepsilon_{\rm vac})b_n.
\]
Para \(j=0\), separar
\(\varepsilon_{\rm vac}\sum_{n\le N}n^{-1}\) da exactamente
\cref{eq:vacancy-gamma}. Una sumaci\'on de Abel acota conservadoramente la
cola posterior a \(N\) por \(2/(N+1)\): un t\'ermino procede del borde y
otro de la variaci\'on total de \(1/n\). La ejecuci\'on exhaustiva declarada usa
\(N=10^8\), enumera las posiciones Beatty de los unos, comprueba los
conjuntos de huecos \(\{21,22\}\) y retornos \(\{6,7\}\), y obtiene como
parte finita
\[
-0.11207107505876366224207105242526397929\ldots.
\]
Al combinar la cota anal\'itica con las guardas de redondeo y frontera
registradas, el certificado finito encierra
\begin{equation}
-0.112071094143613634
<\gamma_{\rm vac}^{+}
<-0.112071055516338732.
\label{eq:vacancy-gamma-interval}
\end{equation}
Este intervalo es un valor terminal con error controlado. No define la
palabra; la palabra y su discrepancia lo preceden.

Una segunda evaluación mediante aritmética multiprecisión y control dirigido
del redondeo proporciona, a dieciocho decimales, el intervalo certificado
\[
-0.112071086058763562
<\gamma_{\rm vac}^{+}
<-0.112071064058763762,
\]
contenido estrictamente en \cref{eq:vacancy-gamma-interval}. El encaje de
ambos intervalos verifica la estabilidad de la parte finita respecto del
truncamiento, de las guardas de frontera y del régimen de precisión. El
resultado depende únicamente de las recurrencias, las cotas de cola y la
aritmética intervalar expuestas en este capítulo.

\section{Contenido operatorio de la representaci\'on helicoidal}

La representaci\'on helicoidal codifica una fase peri\'odica junto con un
incremento de memoria. Posee contenido operatorio cuando existe un cociclo

\[
c(\gamma\eta)=c(\gamma)+c(\eta)
\]

que transporta el acarreo y cuya monodrom\'ia se entrelaza con el operador
aritm\'etico. En ausencia de dicho mapa, la h\'elice constituye s\'olo una
parametrizaci\'on de fase; con el entrelazador, distingue el retorno de fase
del retorno de estado.

\begin{remark}[Control negativo]
La cofinalidad de cilindros no implica conjugaci\'on din\'amica. Un
contraejemplo a la conmutaci\'on entre el desplazamiento TPK y \(T_G\)
impide el entrelazador, pero no afecta a
\cref{eq:levy,eq:gauss-entropy}. En el sector de vacancias, la divergencia
de \cref{eq:vacancy-stieltjes} bajo la cota declarada refuta la familia de
constantes regularizadas.
\end{remark}

\subsection*{Alcance lógico y dominio de validez}
Cofinalidad, identidades de Gauss, L\'evy, entrop\'ia y Lochs: \emph{Teorema interno exacto en el sistema transportado}. GKW y entrelazador din\'amico: \emph{Programa abierto con criterio de refutación}. Densidad, constante y holonom\'ia de vacancia: \emph{Teorema interno exacto con control de discrepancia}.
